\chapter{Testbare Vorhersagen}

\section{Einleitung}

Die IWT und die aus ihr emergierende WDBT+ machen eine Reihe spezifischer, von der Allgemeinen Relativitätstheorie und dem Standardmodell abweichender Vorhersagen.

Diese Vorhersagen sind prinzipiell experimentell überprüfbar – entweder mit bestehenden oder mit zukünftigen Technologien.

Dieses Kapitel fasst die wichtigsten testbaren Vorhersagen zusammen.

\section{Fraktale Skalengesetze in der Kosmologie}

Die fraktale Raumdimension \( D \approx 2,704 \) führt zu gebrochenen Exponenten in kosmologischen Skalengesetzen:

\begin{itemize}
\item \textbf{Galaxien-Dichteprofil:} \( \rho(r) \propto r^{-0,296} \) – flacher als das NFW-Profil.
\item \textbf{Rotationskurven:} \( v(r) \propto r^{0,352} \) – bei sehr großen Radien steigt die Kurve langsam an.
\item \textbf{CMB-Korrelationen:} \( C(\theta) \propto \theta^{-0,296} \).
\item \textbf{Materieentstehung:} \( \dot{\rho}_{\text{cre}}(r) \propto r^{0,296} \).
\end{itemize}

\section{Die maximale Masse kompakter Objekte}

Die Theorie sagt eine natürliche Obergrenze für kompakte Objekte voraus:

\[
M_{\text{BEC}} \approx 3 \times 10^6 M_{\odot}
\]

Beobachtete supermassereiche Objekte mit Massen über dieser Grenze sind keine einzelnen kompakten Objekte, sondern Systeme aus einem BEC-Kern und einer massiven Umgebung.

\section{Gravitationsphysik}

\subsection{Frequenzabhängige Lichtablenkung}

Die ART sagt eine frequenzunabhängige Lichtablenkung voraus. Die WDBT+ sagt dagegen eine frequenzabhängige Ablenkung voraus:

\[
\Delta\phi(\nu) = \Delta\phi_0 \left(1 + \alpha \frac{\nu_0}{\nu}\right) \propto \lambda^2
\]

Dies ist mit hochpräzisen Interferometern testbar.

\subsection{Gravitationswellen-Dispersion}

Die fraktale Raumstruktur führt zu einer frequenzabhängigen Ausbreitungsgeschwindigkeit von Gravitationswellen:

\[
v_{\text{phase}}(\omega) = c \left(1 + \beta \left(\frac{\omega}{\omega_0}\right)^{-0,099} + \dots\right)
\]

Dies ist mit LISA und dem Einstein-Teleskop testbar.

\subsection{Shapiro-Laufzeitverzögerung}

Die ART sagt eine frequenzunabhängige Laufzeitverzögerung voraus. Die WDBT+ sagt eine zusätzliche Frequenzabhängigkeit voraus:

\[
\Delta t_{\text{WG}} \propto \lambda^{-2}
\]

Dies ist mit Pulsar-Timing-Experimenten testbar.

\section{Neutrinophysik}

\subsection{Neutrinomasse}

Die Theorie sagt eine Neutrinomasse von:

\[
m_{\nu,1} \approx 0,1 \, \text{eV}/c^2
\]

mit drei Generationen:
\[
\begin{aligned}
m_{\nu,1} &\approx 0,1 \, \text{eV}/c^2 \\
m_{\nu,2} &\approx 9,2 \times 10^{-4} \, \text{eV}/c^2 \\
m_{\nu,3} &\ll 10^{-4} \, \text{eV}/c^2
\end{aligned}
\]

Dies ist mit KATRIN und dem neutrinolosen Doppelbeta-Zerfall testbar.

\subsection{Neutrino-Oszillationen}

Die fraktale Dispersionsrelation führt zu einer charakteristischen Skalierung der Oszillationslänge:

\[
L_{\text{osc}} \propto E^{0,775}
\]

Das Standardmodell sagt dagegen \( L_{\text{osc}} \propto E \) voraus.

Dies ist mit JUNO, DUNE und Hyper-Kamiokande testbar.

\section{Kosmologie}

\subsection{Die Hubble-Spannung}

Die effektive Hubble-Konstante ist skalenabhängig:

\[
H_{\text{eff}}(d) = H_0 \left( \frac{d}{d_0} \right)^{0,704}
\]

Die Hubble-Spannung ist keine Inkonsistenz, sondern eine Konsequenz der fraktalen Geometrie:

\[
\frac{H_{\text{local}}}{H_{\text{CMB}}} = \left( \frac{d_{\text{local}}}{d_{\text{CMB}}} \right)^{0,704}
\]

\subsection{Stationäres Universum}

Die WDBT+ sagt ein stationäres, nicht-expandierendes Universum voraus.

\begin{itemize}
\item Kein Urknall.
\item Keine Dunkle Materie.
\item Keine Dunkle Energie.
\item Die Rotverschiebung ist kumulative Gravitation, nicht Expansion.
\end{itemize}

\section{Sonnenphysik}

Die WDBT+ sagt voraus, dass der Sonnenwind nicht aus der Sonne selbst stammt, sondern aus neu entstehender Materie.

Die Sonne verliert daher keine signifikante Masse.

Präzise Messungen der Sonnenmasse über Zeiträume von Jahrzehnten sollten diesen Effekt bestätigen – oder die Theorie falsifizieren.

\section{Zusammenfassung der wichtigsten Vorhersagen}

\begin{itemize}
\item Galaxien-Dichteprofil: \( \rho(r) \propto r^{-0,296} \)
\item Rotationskurven: \( v(r) \propto r^{0,352} \)
\item Maximale kompakte Masse: \( M_{\text{BEC}} \approx 3 \times 10^6 M_{\odot} \)
\item Neutrinomasse: \( m_{\nu,1} \approx 0,1 \, \text{eV}/c^2 \)
\item Oszillationslänge: \( L_{\text{osc}} \propto E^{0,775} \)
\item Lichtablenkung: \( \Delta\phi \propto \lambda^2 \)
\item Gravitationswellen-Dispersion: \( v_{\text{phase}} \propto \omega^{-0,099} \)
\item Hubble-Spannung: \( H_{\text{eff}}(d) \propto d^{0,704} \)
\item Stationäres Universum: Keine Expansion, kein Urknall.
\end{itemize}

\section{Fazit}

Die WDBT+ macht spezifische, testbare Vorhersagen, die sie von der ART und dem Standardmodell unterscheiden.

Sollten diese Vorhersagen bestätigt werden, wäre dies ein starkes Indiz für die Richtigkeit der Theorie.

Sollten sie widerlegt werden, wäre die Theorie falsifiziert – ein Zeichen ihrer wissenschaftlichen Qualität.